\section{The \reflector{} Framework}
\label{sec:reflector-framework}

\reflector{} is proposed as a recursive synchronization architecture for AI-assisted software engineering. It is not presented as autonomous general intelligence or as a replacement for engineering judgment. Instead, it turns the prior control primitives into mixed-initiative coordination infrastructure: a systems layer that keeps intent, execution, artifacts, and governance synchronized across repeated delegation cycles \cite{Wiener1948Cybernetics,Shneiderman2020HCAI}.

\begin{figure}[H]
  \centering
  \includegraphics[width=\linewidth]{figure5.png}
  \caption{\reflector{} architecture: human intent and specification anchors are translated into scoped delegation, AI execution, and artifact generation; reflective audits and synchronization checkpoints reconcile outputs with architectural and governance constraints before bounded recursive continuation.}
  \label{fig:figure5}
\end{figure}

\Cref{fig:figure5} illustrates the \reflector{} architecture for this section and captures the directional flow from intent definition to recursive continuation under audit and synchronization control.

\subsection{Architectural Philosophy: Recursive Coherence over Unbounded Automation}

\reflector{} adopts a systems-oriented philosophy: recursive execution should be optimized for coherence, not only for throughput. In practical terms, recursive systems must continuously preserve three alignments: alignment of generated artifacts with current specifications, alignment of local task outcomes with system-level architecture, and alignment of ongoing delegation with human-governed scope.

This philosophy assumes bounded cognition and bounded authority. Human operators cannot continuously reconstruct every intermediate state at AI execution speed, and AI systems should not autonomously revise project intent without explicit governance checkpoints \cite{Simon1957ModelsMan,Amershi2014InteractiveML}. \reflector{} therefore treats synchronization and audit boundaries as first-class architectural requirements rather than optional workflow hygiene.

\subsection{Layered Architecture}

\reflector{} organizes recursive work into five conceptual layers:
\begin{itemize}
  \item \textbf{Human Intent Layer:} defines goals, constraints, priorities, and acceptance conditions.
  \item \textbf{Recursive Execution Layer:} performs decomposition, delegation, generation, and local orchestration.
  \item \textbf{Reflective Audit Layer:} inspects intermediate and final outputs for scope conformance, quality, and drift signals.
  \item \textbf{Synchronization Layer:} reconciles task state, artifact consistency, architectural compatibility, and contextual freshness at explicit checkpoints.
  \item \textbf{Artifact Layer:} stores the inspectable outputs of execution, including code, tests, documentation, plans, figures, and repository state.
\end{itemize}

These layers are conceptually separable but operationally interdependent. Each recursive cycle propagates state across all layers; each checkpoint restores a shared, inspectable model before further delegation.

\begin{table}[H]
  \centering
  \small
  \caption{Core \reflector{} components and their operational outputs.}
  \label{tab:reflector-components}
  \begin{tabularx}{\linewidth}{l X X}
    \toprule
    \textbf{Component} & \textbf{Responsibility} & \textbf{Output} \\
    \midrule
    Intent and scope definition & Encode goals, constraints, and acceptance conditions before delegation. & Scoped task contracts and specification anchors. \\
    Recursive execution engine & Generate and refine implementation artifacts within declared bounds. & Code, documentation, tests, and workflow artifacts. \\
    Reflective audit loop & Evaluate outputs for drift, quality, and requirement conformance. & Audit findings, risk signals, and correction recommendations. \\
    Synchronization checkpoint & Reconcile current outputs with current intent before continuation. & Continue/correct/pause decisions with explicit rationale. \\
    Artifact memory layer & Preserve inspectable state across recursive cycles. & Repository history and reusable synchronization evidence. \\
    \bottomrule
  \end{tabularx}
\end{table}

\subsection{Recursive Orchestration and Scoped Delegation}

Within \reflector{}, recursive orchestration is governed through scoped delegation contracts. A scoped delegation unit defines the task boundary, permitted artifact surface, expected evidence, and required checkpoint conditions before continuation. This converts broad objectives into bounded execution intervals that remain auditable and interruptible.

Scoped delegation is central because recursive systems amplify local assumptions. Without explicit scope boundaries, successful local edits can still increase global incoherence by violating interfaces, invalidating prior rationale, or drifting from current milestone intent. \reflector{} mitigates this amplification by requiring each delegated unit to terminate in an auditable synchronization state before spawning deeper recursion.

\subsection{Reflective Audit and Synchronization Integration}

Reflective auditing and synchronization are integrated as a single control loop. Audits produce evidence about artifact quality, requirement conformance, and emerging drift. Synchronization checkpoints then use that evidence to decide whether the system should continue, correct, rescope, or pause. In this design, audit is the observability function and synchronization is the reconciliation function.

The integration also improves recursive observability. Instead of relying on post hoc interpretation after multiple autonomous steps, \reflector{} introduces recurrent moments where machine-produced artifacts and human interpretation can be re-coupled. This keeps recursive continuation conditional on explicit alignment rather than implicit momentum \cite{Chandy1985DistributedSnapshots,Lamport1978TimeClocks}.

\subsection{Recursive Execution Lifecycle}

\reflector{} models recursive execution as a bounded lifecycle:
\begin{center}
Human Intent $\rightarrow$ Scoped Delegation $\rightarrow$ AI Execution $\rightarrow$ Artifact Generation $\rightarrow$ Reflective Audit $\rightarrow$ Synchronization Checkpoint $\rightarrow$ Alignment Correction $\rightarrow$ Recursive Continuation
\end{center}

The lifecycle is intentionally cyclic rather than linear. Each continuation step is re-authorized based on current synchronized state, not assumed from prior progress. This supports iterative stabilization: recursive work can accelerate when alignment is strong and contract when drift signals increase.

\subsection{Specification and Artifact Synchronization Anchors}

\reflector{} treats specifications and repository artifacts as synchronization anchors. Issues and task descriptions define scoped intent boundaries; pull requests and repository state provide inspectable evidence of execution; documentation and generated diagrams preserve architectural rationale for subsequent recursive cycles. Together, these artifacts function as recursive memory: they carry forward context while remaining open to audit and correction.

This perspective avoids over-specifying implementation while preserving operational plausibility. \reflector{} does not require a single toolchain; it requires that whatever toolchain is used exposes inspectable artifacts and explicit checkpoint semantics that keep recursive work governable.

\subsection{Human-AI Mixed-Initiative Governance and Transition}

\reflector{} is a mixed-initiative governance model in which humans retain directional and normative authority while AI systems provide high-throughput execution within bounded scope. Synchronization checkpoints maintain coherence between these roles by ensuring that continuation remains a governed decision rather than an automatic consequence of local completion.

With this architecture in place, the next section can move from conceptual framing to operational workflow. We therefore transition from architectural definition to practical demonstration of recursive issue orchestration, artifact synchronization, and checkpoint-driven continuation in realistic repository contexts (\Cref{sec:operational-demonstration}).
